\originalchapter{从谱线位移到均匀宇宙的初始值问题}
本章对应8.286第3--8讲的观测和牛顿动力学入口，补充后面的相对论背景章。课程从可测量的频率开始，再问同一种膨胀律如何被不同位置的观测者共同描述。狭义相对论的一般原理见广义相对论卷；这里保留理解谱线和宇宙学坐标必需的推导。距离、速度和时间的定义必须先统一，才有可能比较数据和模型。

\originalsection{声波的Doppler效应与光的Doppler效应}
先在静止介质中考虑一维声波。波速为$u$，源周期为$T_s$；源向观察者运动，速率$v_s$。两次发射之间源前进$v_sT_s$，前一波峰却前进$uT_s$，所以前方波长为$\lambda=(u-v_s)T_s$。若观察者也向源运动，速率$v_o$，其遇到波峰的速率为
\[
 f_o=f_s\frac{u+v_o}{u-v_s}.
\]
这一区分源和观察者的公式依赖\emph{介质静止系}。使源远离观察者取$v_s<0$，观察者远离源取$v_o<0$。它不能直接令$u=c$后用于真空光波，因为光没有这一特殊介质系。

对相距渐增的光源与观察者，取观察者惯性系，光源退行速率$v=\beta c$，源的固有发射周期为$T_s$。该系里发射间隔是$\gamma T_s$，其中$\gamma=(1-\beta^2)^{-1/2}$；第二个波峰还需多走$v\gamma T_s$。故接收间隔为
\[
 T_o=\gamma(1+\beta)T_s,\qquad
 1+z=\frac{T_o}{T_s}=\sqrt{\frac{1+\beta}{1-\beta}}.
\]
这既含传播路程差，也含钟速差。反解$\beta=[(1+z)^2-1]/[(1+z)^2+1]$，小$z$时才有$v\simeq cz$。一般方向最清楚的定义是局部四动量测频：$E_{\rm obs}=-p_\mu u^\mu$，$\nu=E_{\rm obs}/h$，这样方向和参考系不会混淆。对于观测系中速度与光传播方向夹角$\theta$，有$\nu_o=\nu_s/[\gamma(1-\beta\cos\theta)]$。横向$\cos\theta=0$仍有红移$1+z=\gamma$。

两事件若在物体自己的静止系同地发生，Lorentz变换给$\Delta t=\gamma\Delta\tau$。测运动尺长却需在\emph{测量者系同时}取两端，才得$L=L_0/\gamma$；两个尺端在自身静止系同时的读数，变换后一般不同时。课程第3讲的分步波峰图正是在提醒这一区别。宇宙学红移来自沿途不同局部共动系的累计传播，不能把远处退行率$HD$无条件当作上述一个惯性系的$v$。

\originalsection{角度、视差与距离梯级}
周年视差$p$是地球轨道一AU基线在天体处张开的半角。小角关系$D=\mathrm{AU}/p$中$p$须用弧度；$p=1$角秒定义一秒差距，因此
\[
 D(\mathrm{pc})=\frac1{p(\mathrm{arcsec})},\qquad
 1\ \mathrm{pc}=206265\ \mathrm{AU}\simeq3.086\times10^{16}\ \mathrm m.
\]
这是几何测距，不需先假设恒星光度。再用已知距离的恒星校准光度或周期关系，才能推广到造父变星和超新星。若一个校准层的零点有偏移，它会传递到更远的层；单一光通量本身不能区分远而亮、近而暗。

谱线红移通过静止实验室波长比较得$z=(\lambda_o-\lambda_s)/\lambda_s$。近邻星系的红移与独立距离呈近线性关系，这提出均匀膨胀的模型。8.286第4讲分别讨论Lemaître于1927年的理论联系与Hubble于1929年的观测图；本卷保留这些年代，原历史照片和曲线不复制。现代拟合还需扣除本动速度；例如300 km/s的本动速度在1000 km/s退行速度上占30\%，在30000 km/s上占1\%。这说明近邻更利于几何校准，却不一定最适合单独估宇宙膨胀率。

\originalsection{均匀膨胀为什么没有特殊中心}
在平直近似下，共动网格位置记为$\mathbf x$，物理位置$\mathbf r=a(t)\mathbf x$。从原点O换到共动位置$\mathbf x_P$，相对坐标是$\mathbf r'=a(\mathbf x-\mathbf x_P)$，于是$\dot{\mathbf r}'=H\mathbf r'$。每个共动观察者都看到同样的Hubble律。气球表面的点是二维类比：表面上没有中心点，而球心属于额外的嵌入空间，并不对应表面上的观测者。三维宇宙无需假定真实存在一个外部空间。

尺度因子是网格间距的共同倍数；固定网格只是描述平均流动，不保证真实星系严格静止于网格。原课程Guth用$a$的单位“米／格”，$r$的单位“格”，曲率参数$k$的单位“格$^{-2}$”。本卷使用无量纲$a$和有长度的$\chi$。令本卷的$a=a_G/a_{G0}$、$\chi=a_{G0}r_G$，就有$K=k_G/a_{G0}^2$；两种写法描述同一几何。不能把其中一种的单位要求误报成另一种的错误。

\originalsection{无限均匀质量的引力为何需要边界条件}
牛顿理论对有限分布可由势积分定义引力。无限均匀静态分布则有边界问题：若将越来越大的球总取以P为中心，壳层定理会给P处$\mathbf g=0$；若总以Q为中心，P位于均匀球内部，得到$\mathbf g=-4\pi G\rho(\mathbf r_P-\mathbf r_Q)/3$。不同极限选择了不同的加速度零点。没有声明极限和边界条件，不能由“各方向相同”断言所有位置的引力同时为零。

Poisson方程给$\nabla\cdot\mathbf g=-4\pi G\rho$，处处为零的场在$\rho>0$时不满足该式。均匀模型在有限邻域内的一个解是$\mathbf g=-4\pi G\rho\mathbf r/3$，加上任意随时间变化的均匀加速度仍不改变散度。两个质点的\emph{相对}加速度
\[
 \ddot{\mathbf r}_{12}=-\frac{4\pi G\rho}{3}\mathbf r_{12}
\]
却不依赖这一共同加速度。第5--6讲用Newton与Bentley通信引出的难点，应理解为无限静态极限与动力学边界条件的问题；物理上可测的相对运动是均匀宇宙的可靠入口。

\originalsection{有限尘埃球怎样给出可求解的模型}
取初始半径内密度$\rho_i$均匀、初始流速$H_i\mathbf r_i$，忽略压强与外界潮汐。壳层不交叉时，$r(t)=a(t)r_i$，取$a_i=1$；包围质量$M=4\pi\rho_i r_i^3/3$恒定，外壳不影响其牛顿引力。故
\[
 \ddot a=-\frac{4\pi G\rho_i}{3a^2},\qquad
 \frac{\dot a^2}{2}-\frac{4\pi G\rho_i}{3a}=E_a.
\]
积分常数$E_a=H_i^2/2-4\pi G\rho_i/3$有$\mathrm{s^{-2}}$量纲。$E_a<0$存在最大尺度，$E_a=0$给$a\propto t^{2/3}$，$E_a>0$在晚期近线性膨胀。这些结论限定为无压强、无宇宙学常数的模型；加入真空能后空间正曲率不再保证再坍缩。

牛顿球是初始值模型和相对运动的解释工具。它没有推出整个宇宙的空间曲率、光线传播或压力引力作用。下面先建立常曲率的实际距离，再在相对论背景章明确这些缺口如何由度规和Einstein方程补上。
